% Factor-level analysis of the two LLM-evolved alpha books.
% Build: tectonic factor_book_analysis.tex   (figures expected in figures_en/)
\documentclass[11pt,a4paper]{article}

\usepackage[margin=2.6cm]{geometry}
\usepackage{graphicx}
\usepackage{booktabs}
\usepackage{makecell}
\usepackage{float}
\usepackage{caption}
\usepackage{enumitem}
\usepackage{xcolor}
\usepackage{microtype}
\usepackage{newtxtext,newtxmath}
\usepackage[hidelinks,pdfauthor={Luca Wuerker},
            pdftitle={Factor-Level Analysis of Two LLM-Evolved Alpha Books}]{hyperref}

\definecolor{grey}{RGB}{96,93,84}
\captionsetup{font=small,labelfont=bf,width=.92\textwidth}
\setlist{itemsep=2pt,topsep=4pt}
\newcommand{\fid}[1]{{\small\ttfamily #1}}

\title{\bfseries Factor-Level Analysis of Two LLM-Evolved Alpha Books}
\author{Luca W\"urker}
\date{August 2026}

\begin{document}
\maketitle

\begin{abstract}
\noindent Factor-level comparison of the alpha books published by two evolutionary
factor-research runs --- GPT-5.6 Terra at the full L4 configuration and Claude
Opus~5 --- benchmarked against the 101 formulaic alphas of Kakushadze (2016) on a
survivorship-bias-free Nasdaq-100 point-in-time panel (2010--2026). Signals are
evaluated directly: information coefficients per factor and per combined book on
three disjoint windows, correlation structure, effective diversity, and novelty
relative to the formulaic library. No strategy construction, position sizing, or
costs enter this analysis.
\end{abstract}

\section{What is measured}

Each evolution run publishes a \emph{book} of factors: small programs mapping the
market panel (prices, volumes, fundamentals, analyst estimates) to one signal value
per (day, stock). Two kinds of information coefficient (IC) are reported, both
against 6-day forward returns ($h=6$ bars, the horizon both runs searched at):

\begin{itemize}
\item \textbf{Per-factor IC.} For each factor and each stock, the Pearson
correlation between the factor's signal series and that stock's own 6-day
forward-return series; the observation-weighted average across stocks is the
factor's \emph{pooled per-underlying IC}. This is the statistic family the
evolutionary harness itself optimises, so it is the honest yardstick for the
search.
\item \textbf{Combined-book IC.} All signals of a book are standardised per stock
and fed to a LightGBM regressor (the same model class the runs used to score
marginal value), fitted once on the DEV window only, predicting the 6-day forward
return. The pooled per-underlying IC of that single combined prediction measures
the book as a whole, including nonlinear interactions between factors.
\end{itemize}

Three evaluation windows follow the runs' own 60/20/20 split convention of the
research panel (2010\,$\to$\,2024-07):

\begin{table}[H]
\centering\small
\begin{tabular}{@{}llrp{6.4cm}@{}}
\toprule
\textbf{Window} & \textbf{Period} & \textbf{Bars} & \textbf{Role} \\
\midrule
DEV (IS+VAL) & 2010-01-04 -- 2021-08-25 & 2\,932 &
Search window: factors were selected on (parts of) this data; the combined model
is fitted here. Combined ICs on DEV are in-sample. \\
TEST & 2021-08-26 -- 2024-07-26 & 733 &
Final 20\,\% of the research panel, never revealed to any generation of the
search. \\
FORWARD & 2024-07-29 -- 2026-07-27 & 500 &
Entirely outside the research panel (data obtained after the runs finished); the
cleanest out-of-sample evidence. \\
\bottomrule
\end{tabular}
\end{table}

\noindent\textit{\small\color{grey}Comparability caveat: no Opus~5 run exists at
the L4 configuration (the Opus ladder stopped after L2 when its budget was
exhausted). The comparison is therefore L4\_terra\_s0 (GPT-5.6, GraphRAG
retrieval, 8 mechanism groups, 44 published factors, \$87) vs.\ L2\_opus5\_s0
(Claude Opus~5, no retrieval, 1 group $\times$ 4 demes, 18 factors, \$452). Model
and configuration are confounded; differences cannot be attributed to the model
alone. Four of the 44 Terra factors produce degenerate (variance-free) signals on
the extended analysis panel and are excluded from per-factor and correlation
statistics; they remain in the combined fits, where they contribute nothing.}

\section{Combined books: IC by window}

\begin{table}[H]
\centering\small
\begin{tabular}{@{}lrrr@{}}
\toprule
\textbf{Book} & \textbf{DEV (in-sample)} & \textbf{TEST} & \textbf{FORWARD} \\
\midrule
101 formulaic alphas alone  & 0.210 & 0.026 & 0.024 \\
Terra book (44)             & 0.213 & 0.024 & 0.033 \\
Opus book (18)              & 0.158 & \textbf{0.045} & 0.027 \\
Terra + Opus (62)           & 0.221 & 0.041 & 0.039 \\
Terra + 101 (145)           & 0.236 & 0.031 & \textbf{0.041} \\
Opus + 101 (119)            & 0.216 & 0.039 & 0.027 \\
Terra + Opus + 101 (163)    & 0.244 & 0.039 & 0.040 \\
\bottomrule
\end{tabular}
\caption{Pooled per-underlying IC of the combined LightGBM signal ($h=6$), fitted
on DEV only. DEV values are in-sample and shown only to gauge the generalisation
gap.}
\end{table}

\begin{figure}[H]
\centering
\includegraphics[width=\textwidth]{figures_en/fig3_combined_ic.png}
\caption{Combined-book IC per window (LightGBM fitted on DEV only).}
\end{figure}

\begin{itemize}
\item \textbf{Opus generalises best on its own.} TEST IC 0.045 from a DEV fit of
0.158 --- a retention of 28\,\%, vs.\ 11\,\% for Terra and 12\,\% for the
formulaic library.
\item \textbf{Terra is a marginal-value book.} Alone it sits at library level
(TEST 0.024), but added to the 101 alphas it lifts the FORWARD IC from 0.024 to
0.041 --- consistent with the L4 run's selection objective (leave-one-out marginal
contribution against exactly this reference book).
\item \textbf{The two books combine well.} Terra + Opus reaches TEST 0.041 /
FORWARD 0.039 from only 62 factors; adding the 101 alphas on top changes little
(0.039 / 0.040).
\end{itemize}

\section{Per-factor ICs: distribution and degradation}

\begin{table}[H]
\centering\small
\begin{tabular}{@{}lrrrrrr@{}}
\toprule
\textbf{Book} & $n$ & \makecell*[r]{mean $|$IC$|$\\DEV} &
\makecell*[r]{mean $|$IC$|$\\TEST} & \makecell*[r]{mean $|$IC$|$\\FWD} &
\makecell*[r]{sign ret.\\DEV$\to$TEST} & \makecell*[r]{corr\\(DEV, TEST)} \\
\midrule
Terra (L4)          & 40  & 0.0082 & 0.0106 & 0.0133 & 80\,\% & 0.55 \\
Opus (L2-evolution) & 18  & 0.0135 & 0.0176 & 0.0152 & 83\,\% & 0.65 \\
101 alphas          & 101 & 0.0143 & 0.0110 & 0.0126 & 75\,\% & 0.66 \\
\bottomrule
\end{tabular}
\caption{Per-factor pooled IC statistics by window.}
\end{table}

\begin{figure}[H]
\centering
\includegraphics[width=\textwidth]{figures_en/fig1_ic_scatter.png}
\caption{Each point is one factor: DEV IC against TEST IC (left) and FORWARD IC
(right). Points do not collapse below the diagonal: average per-factor ICs do not
degrade out of sample, because the search never selected on standalone IC (it
selected on marginal contribution to the book), so standalone IC was never
over-optimised. Sign retention is 80--89\,\% for the LLM books vs.\ 74--75\,\%
for the formulaic library.}
\end{figure}

\paragraph{Key finding: the apparent overfitting sits in the combining model, not
in the factors.} The combined-book IC drops by roughly an order of magnitude from
DEV to TEST (e.g.\ Terra 0.213\,$\to$\,0.024), which at first glance looks like
heavy factor overfitting. The per-factor evidence contradicts that reading:
individual factor ICs are as strong on TEST and FORWARD as on DEV (mean $|$IC$|$
flat to rising, sign retention 80--89\,\%). What collapses out of sample is the
fitted LightGBM \emph{combination} --- its in-sample IC reflects the model
memorising DEV-specific interaction patterns, an in-sample artefact that says
little about the factors themselves. The honest generalisation gap of the books
is the much smaller one visible per factor, plus the 2021 regime break documented
in Section~\ref{sec:regime}.

\section{Strongest individual factors}

\begin{table}[H]
\centering\footnotesize
\begin{tabular}{@{}llrrr@{}}
\toprule
\textbf{Factor (Opus, top 8 by $|$TEST IC$|$)} & \textbf{Category} &
\textbf{DEV} & \textbf{TEST} & \textbf{FWD} \\
\midrule
\fid{comparable\_basket\_repair\_rate\_activity}   & stat.\ arbitrage & 0.024 & \textbf{0.062} & 0.024 \\
\fid{rice\_consistency\_wedge\_gated\_crossing}    & mean reversion   & 0.016 & 0.041 & 0.019 \\
\fid{run\_surprisal\_vs\_flip\_baseline\_exhaust.} & mean reversion   & 0.020 & 0.037 & 0.024 \\
\fid{churn\_resilience\_gated\_sweep\_debt\_snap.} & microstructure   & 0.040 & 0.031 & 0.021 \\
\fid{reignition\_hazard\_pending\_spell\_vs\_paid} & microstructure   & 0.037 & 0.030 & 0.033 \\
\fid{program\_gated\_elasticity\_factor\_return}   & stat.\ arbitrage & 0.002 & 0.019 & 0.024 \\
\fid{absorbed\_wick\_routed\_unowed\_transient}    & microstructure   & 0.033 & 0.017 & 0.029 \\
\fid{volume\_evenness\_entropy\_innovation}        & microstructure   & $-$0.001 & 0.016 & $-$0.000 \\
\bottomrule
\end{tabular}
\end{table}

\begin{table}[H]
\centering\footnotesize
\begin{tabular}{@{}llrrr@{}}
\toprule
\textbf{Factor (Terra, top 8 by $|$TEST IC$|$)} & \textbf{Category} &
\textbf{DEV} & \textbf{TEST} & \textbf{FWD} \\
\midrule
\fid{nonlinear\_downside\_coskew\_compensation}    & volatility     & 0.004 & 0.034 & 0.006 \\
\fid{peer\_hawkes\_contagion\_underreaction}       & microstructure & 0.036 & 0.029 & 0.037 \\
\fid{consumption\_tail\_amplification\_premium}    & other          & 0.012 & 0.027 & 0.023 \\
\fid{news\_synchrony\_forcedflow\_conversion}      & microstructure & $-$0.001 & $-$0.026 & 0.014 \\
\fid{stale\_surprise\_fear\_state\_cascade\_diff.} & sentiment      & 0.007 & 0.023 & 0.026 \\
\fid{earnings\_turnover\_impedance\_resolution}    & microstructure & 0.019 & 0.022 & $-$0.002 \\
\fid{coherent\_surprise\_response\_gap\_drift}     & sentiment      & $-$0.010 & 0.021 & 0.002 \\
\fid{tail\_cluster\_resilient\_liquidity\_release} & mean reversion & 0.026 & 0.017 & 0.026 \\
\bottomrule
\end{tabular}
\caption{Strongest factors of each book, ranked by $|$TEST IC$|$ (pooled
per-underlying, $h=6$).}
\end{table}

Reference: the best formulaic alphas on TEST are \fid{alpha\_024} (0.052) and
\fid{alpha\_035} (0.037). The single best factor of the whole analysis is Opus's
\fid{comparable\_basket\_repair\_rate\_activity} (TEST 0.062 from a DEV of only
0.024).

\begin{figure}[H]
\centering
\includegraphics[width=\textwidth]{figures_en/fig2_top_factors_opus.png}
\caption{Opus: top 15 factors by $|$DEV IC$|$, with TEST and FORWARD alongside.}
\end{figure}

\begin{figure}[H]
\centering
\includegraphics[width=\textwidth]{figures_en/fig2_top_factors_terra.png}
\caption{Terra: top 15 factors by $|$DEV IC$|$.}
\end{figure}

\section{Correlation structure and diversity}

Signal correlations are computed on cross-sectionally standardised signals over
the DEV window. Within-book redundancy is low for both LLM books (mean $|\rho|$:
Terra 0.093, Opus 0.094, library 0.120), but \emph{efficiency} differs: the
participation ratio of the correlation spectrum gives Opus 12.7 effective factors
out of 18 (71\,\%), Terra 19.7 / 41 (48\,\%), and the formulaic library only
25.8 / 101 (26\,\%).

The cross-blocks are weak: mean $|\rho|$ Terra$\times$Opus 0.056 (max 0.62),
Terra$\times$library 0.076, Opus$\times$library 0.093. The two models explored
largely orthogonal factor spaces: Opus concentrated on microstructure (12 of 18
factors), Terra spread across sentiment, volatility, fundamentals and
microstructure. Pooling everything yields 35.0 effective factors from 160.

\begin{figure}[H]
\centering
\includegraphics[width=\textwidth]{figures_en/fig4_corr_heatmaps.png}
\caption{Within-book correlations, hierarchically ordered. Terra carries one
visible redundancy cluster (liquidity/downside family, pairs up to
$|\rho|\approx 0.85$); Opus a small four-factor wick/absorption block.}
\end{figure}

\begin{figure}[H]
\centering
\includegraphics[width=.8\textwidth]{figures_en/fig5_corr_blocks.png}
\caption{Full correlation matrix with block boundaries. The library block
(bottom right) is visibly denser; the off-diagonal blocks stay pale.}
\end{figure}

\begin{figure}[H]
\centering
\includegraphics[width=.85\textwidth]{figures_en/fig7_effective_n.png}
\caption{Nominal vs.\ effective factor counts (participation ratio of the
correlation eigenvalues).}
\end{figure}

\section{Novelty relative to the formulaic library}

For each evolved factor, the maximum $|$correlation$|$ to any of the 101
formulaic alphas measures how much of it the library already spans. The median
is 0.28 (Terra) and 0.29 (Opus): the typical evolved factor is not a disguised
formulaic alpha. The most library-like factors of both books attach to the same
two alphas (\fid{alpha\_033}, \fid{alpha\_038}), with a maximum of 0.78.

\begin{table}[H]
\centering\footnotesize
\begin{tabular}{@{}lllrr@{}}
\toprule
\textbf{Book} & \textbf{Factor} & \textbf{Nearest alpha} &
\makecell*[r]{max $|\rho|$\\library} & \makecell*[r]{max $|\rho|$\\own book} \\
\midrule
Terra & \fid{idiosyncratic\_downside\_hawkes\_branch}  & \fid{alpha\_033} & 0.78 & 0.63 \\
Opus  & \fid{reignition\_hazard\_pending\_spell\_vs\_paid} & \fid{alpha\_038} & 0.65 & 0.54 \\
Terra & \fid{conditional\_impact\_copula\_information} & \fid{alpha\_033} & 0.63 & 0.73 \\
Opus  & \fid{tsrv\_kalman\_gain\_transient\_deviation} & \fid{alpha\_033} & 0.60 & 0.41 \\
Terra & \fid{tail\_cluster\_resilient\_liquidity\_release} & \fid{alpha\_038} & 0.59 & 0.36 \\
Terra & \fid{sparse\_liquidity\_vacuum\_repair}        & \fid{alpha\_062} & 0.59 & 0.85 \\
\bottomrule
\end{tabular}
\caption{The six most library-like evolved factors.}
\end{table}

\begin{figure}[H]
\centering
\includegraphics[width=\textwidth]{figures_en/fig6_zoo_novelty.png}
\caption{Left: distribution of library-similarity (binned in 0.1-wide buckets);
Terra has a solid block of genuinely novel factors below 0.15. Right: factors
above the diagonal are more redundant within their own book than with the
library --- Terra's redundancy is home-grown (the cluster of Figure~5), not
library imitation.}
\end{figure}

\section{IC through time: the 2021 regime break}\label{sec:regime}

The pooled per-underlying IC is one number per window and has no daily analogue,
so the time-resolved view uses the daily \emph{cross-sectional} IC (each day, the
correlation across stocks between signal and forward return) as a diagnostic; its
level is not comparable to the pooled numbers above.

\begin{figure}[H]
\centering
\includegraphics[width=\textwidth]{figures_en/fig8_rolling_ic.png}
\caption{63-day rolling mean of the daily cross-sectional IC of the combined
signals (fitted on DEV only). All three books --- including the formulaic
library, which was never fitted to this panel --- collapse simultaneously at the
2021 TEST boundary: the cross-sectional factor regime in the Nasdaq-100 itself
changed (mega-cap concentration), so the DEV$\to$TEST drop is not pure
overfitting. The Opus signal stays above zero most often out of sample; the
library turns clearly negative in 2025--26.}
\end{figure}

\section{Conclusions}

\begin{itemize}
\item \textbf{Per-factor quality: Opus.} Individually stronger and more stable
factors (mean $|$TEST IC$|$ 0.018 vs.\ 0.011), the best standalone
generalisation of the combined book (TEST 0.045), and the highest
effective-dimension efficiency --- at $5\times$ the cost (\$452 vs.\ \$87) and
with only 18 factors.
\item \textbf{Breadth and marginal value: Terra.} More factors, more genuine
novelty, the largest diversity gain to the pooled book, and the best FORWARD
results appear only in combination (Terra + 101: 0.041) --- exactly the
selection objective of the L4 configuration.
\item \textbf{Near-orthogonal books.} Mean cross-correlation of 0.056 between
the two books; combining them (with or without the formulaic library) is the
natural next step and already outperforms either book alone on TEST.
\item \textbf{The DEV$\to$TEST gap is a model artefact, not factor
overfitting.} Per-factor ICs hold or improve out of sample while only the
fitted combination collapses; in-sample combined ICs should never be quoted as
evidence of book quality.
\item \textbf{Attribution caveat.} Model (GPT-5.6 vs.\ Opus~5) and
configuration (L4 GraphRAG / 8 groups vs.\ L2 no-retrieval / 1 group) are
confounded; a clean model comparison would require the never-run L4 Opus arm.
\end{itemize}

\end{document}
